\documentclass{article}
\usepackage[UTF8]{ctex}
\usepackage{amsmath,amssymb,amsfonts}
\usepackage{enumitem}
\usepackage{geometry}
\geometry{a4paper, margin=1in}
% === 错题本排版优化（不参与题目标识） ===
\setlength{\parindent}{0pt}
\setlength{\parskip}{0.35em plus 0.10em minus 0.08em}
\setlength{\emergencystretch}{3em}
\setlist{itemsep=0.25em, topsep=0.40em, parsep=0pt}
\allowdisplaybreaks[2]
% === 错题本排版优化结束 ===
\begin{document}
\section*{武汉大学2024年数学分析真题}
\section*{130．武汉大学2024年数学分析试卷}

\subsection{题目1}
1． \(\lim _{x \rightarrow 0} \frac{\sin x-\arctan x}{\tan x-\arcsin x}\) ．

\subsection{题目2}
2．\(y=\frac{\arcsin x}{\sqrt{1-x^{2}}}\) ，求 \(y^{(2023)}(0)\) ．

\subsection{题目3}
3． \(\lim _{n \rightarrow \infty} \sum_{k=1}^{n} \sin \frac{k}{n} \cdot \sin \frac{k}{n^{2}}\) ．

\subsection{题目4}
1．以 \(x=t y\) 参数化曲线 \(x^{2}+y^{3}=x y\) ，求曲线所围区域的面积．

\subsection{题目5}
2．求曲面积分

\[
\iint_{S}(x y+y z+z x) \mathrm{d} S
\]

其中 \(S\) 是 \(z=\sqrt{x^{2}+y^{2}}\) 被 \(x^{2}+y^{2}=2 a y\) 截掉的部分。

\subsection{题目6}
3．用含参量积分求 \(\int_{0}^{1} \frac{\arctan x}{x \sqrt{1-x^{2}}} \mathrm{~d} x\) ．

\begin{itemize}

\subsection{题目7}
\end{itemize}

\subsection{题目1}
原题边界“1．”：证明：
\[
\frac{2 n}{2 n+1} \leq \frac{\left(1+\frac{1}{n}\right)^{n}}{e} \leq \frac{2 n+1}{2 n+2}(n \geq 1) .
\]

\subsection{题目8}

\subsection{题目2}
原题边界“2．”：设 \(f(x)\) 在 \((0,+\infty)\) 上三次可导，且 \(\lim _{x \rightarrow+\infty} f(x)\) 与 \(\lim _{x \rightarrow+\infty} f^{\prime \prime \prime}(x)\) 均存在，证明：
\[
\lim _{x \rightarrow+\infty} f^{\prime}(x)=\lim _{x \rightarrow+\infty} f^{\prime \prime}(x)=\lim _{x \rightarrow+\infty} f^{\prime \prime \prime}(x)=0 .
\]




\subsection{题目9}

\subsection{题目3}
原题边界“3．”：设 \(f(x)\) 是 \([-1,1]\) 上的连续函数，证明：
\[
\lim _{\epsilon \rightarrow 0^{+}} \int_{-1}^{1} \frac{\varepsilon f(x)}{\varepsilon^{2}+x^{2}} \mathrm{~d} x=\pi f(0)
\]

\subsection{题目10}

\subsection{题目4}
原题边界“4．”：设 \(x>0, y>0\) ，证明：
\[
x y-e^{x-1} \leq y \ln y
\]

\subsection{题目11}
5．证明： \(\int_{0}^{1}(x \ln x)^{n} \mathrm{~d} x=\frac{(-1)^{n} n!}{(n+1)^{n+1}}\) ，并利用此结论证明：

\[
\int_{0}^{1} x^{-x} \mathrm{~d} x=\sum_{n=1}^{\infty} \frac{1}{n^{n}}
\]

\subsection{题目12}
6．设 \(\left\{f_{n}(x)\right\}\) 是区间 \([a, b]\) 上一致收敛于 \(f\) 的可积函数列．证明：\(f\) 在 \([a, b]\) 上可积，且

\[
\int_{a}^{b} f(x) \mathrm{d} x=\lim _{n \rightarrow \infty} \int_{a}^{b} f_{n}(x) \mathrm{d} x .
\]
\end{document}
